% ===========================================================================
\chapter{Matching}\label{ch:matching}

Everything the engine does comes down to one question, asked over and over:
does this pattern match this fact? \code{follows(X, bo)} matches
\code{follows(ann, bo)} if \code{X} is \code{ann}. It does not match
\code{follows(ann, cy)} at all. \code{follows(X, X)} matches only a fact
whose two arguments are the same.

\section{Bindings}

The answer to ``does it match'' is either ``no'' or ``yes, if \code{X} is
\code{ann}''. The second kind of answer gives values to variables, and a
set of such values is called \emph{bindings}.

Our bindings never change once they are made. Binding one more variable
makes a new \code{Bindings} that copies the old values and adds one. That
costs a copy, but it pays for itself in the next chapter, where a rule's
body is matched one condition at a time, and each condition may extend the
same bindings in several different ways. If the bindings could change, each
of those ways would have to undo its changes before the next one tried.

\begin{newfile}{bindings.nex}
-- bindings.nex: giving variables values, and matching atoms against facts.

intern Terms

-- Values for some variables. A Bindings never changes once made: binding
-- one more variable makes a new one, so a search can try several ways
-- forward from the same bindings.
class Bindings
create
  empty
  do
    values := {}
  end

  extending(b: Bindings, name: String, t: Term)
  require
    unbound: not b.has(name)
    constant: not is_variable(t)
  do
    values := b.values.clone
    values.put(name, t)
  end
feature
  values: Map[String, Term]

  has(name: String): Boolean
  do
    result := values.contains_key(name)
  end

  value_of(name: String): Term
  require
    bound: has(name)
  do
    result := values.get(name)
  end

  -- "X = ann, Y = bo", in name order.
  text: String
  do
    result := ""
    across values.keys.sort as name do
      if result /= "" then
        result := result + ", "
      end
      result := result + name + " = " + term_text(values.get(name))
    end
  end
end
\end{newfile}

The second constructor's preconditions say what binding means here: a
variable is bound once, and to a constant, never to another variable.

\section{Substitution}

Applying bindings to an atom replaces each bound variable by its value:
\code{follows(X, Y)} under \code{X = ann} becomes \code{follows(ann, Y)}.

\begin{continued}{bindings.nex}
-- `t` with its variable replaced by its value, if it has one.
function substitute_term(t: Term, b: Bindings): Term
do
  result := t
  let name := variable_name(t)
  if name /= "" and b.has(name) then
    result := b.value_of(name)
  end
end

-- `a` with every bound variable replaced by its value.
function substitute(a: Atom, b: Bindings): Atom
do
  result := create Atom.make(a.predicate, a.args.map(fn (t: Term): Term do
    result := substitute_term(t, b)
  end))
ensure
  same_predicate: result.predicate = a.predicate
end
\end{continued}

\section{Matching an atom against a fact}

Matching goes argument by argument, carrying the bindings along. A constant
in the pattern must equal the fact's constant. A variable that is already
bound must have the value the fact has there. A variable that is not yet
bound becomes bound to it.

\begin{continued}{bindings.nex}
-- The bindings, extending `b`, that make `pattern` equal to `fact`; or
-- nil if there are none.
function match_atom(pattern: Atom, fact: Atom, b: Bindings): ?Bindings
require
  fact_is_ground: fact.is_ground
do
  result := nil
  if pattern.predicate = fact.predicate and pattern.arity = fact.arity then
    let current := b
    let failed := false
    from let i := 0 until i = pattern.arity or failed do
      let next := match_term(pattern.args.get(i), fact.args.get(i), current)
      if next = nil then
        failed := true
      else
        current := next
      end
      i := i + 1
    end
    if not failed then
      result := current
    end
  end
ensure
  matched: is_match(pattern, fact, result)
end

-- Is `b` nil, or does it make `pattern` equal to `fact`?
function is_match(pattern, fact: Atom, b: ?Bindings): Boolean
do
  result := true
  if b /= nil then
    result := substitute(pattern, b).text = fact.text
  end
end

-- The bindings, extending `b`, that make `t` equal to the constant `c`.
function match_term(t: Term, c: Term, b: Bindings): ?Bindings
require
  constant: not is_variable(c)
do
  let name := variable_name(t)
  if name = "" then
    result := when t = c then b else nil end
  elseif b.has(name) then
    result := when b.value_of(name) = c then b else nil end
  else
    result := create Bindings.extending(b, name, c)
  end
end
\end{continued}

The postcondition \code{matched} says what a match \emph{means}: if
\code{match_atom} returns bindings, then substituting them into the pattern
gives exactly the fact. It is a short sentence that is easy to get wrong in
code (forgetting the repeated-variable case, for example), and with it
every match the engine ever makes is checked. The check needs an
\code{if} to look inside a \code{?Bindings}, so it lives in its own
query, \code{is_match}.

\section{Tests}

Rules can also test their variables: \code{X != Y}, \code{Km < 1000}. By
the time a test is checked, its variables have values, so a test compares
two constants. Equality works on any constants; the order comparisons only
make sense for numbers, so they need two integers and are false otherwise.

\begin{continued}{bindings.nex}
-- Is `left op right` true? Equality works on any constants; the order
-- comparisons need two integers.
function holds(op: String, left, right: Term): Boolean
require
  constants: not is_variable(left) and not is_variable(right)
  known: ["=", "!=", "<", "<=", ">", ">="].contains(op)
do
  if op = "=" then
    result := left = right
  elseif op = "!=" then
    result := left /= right
  else
    result := false
    match left of
      Int(value as a) then
        match right of
          Int(value as b) then
            result := (op = "<" and a < b) or (op = "<=" and a <= b)
                      or (op = ">" and a > b) or (op = ">=" and a >= b)
          _ then result := false
        end
      _ then result := false
    end
  end
end
\end{continued}

\section{Trying it out}

\begin{tryit}{try_matching.nex}
intern Parser
intern Bindings

-- The one atom in `text`.
function atom_of(text: String): Atom
do
  match parse_program("?- " + text + ".") of
    Ok as ok then result := ok.value.queries.get(0)
    Err as e then raise e.error
  end
end

function try_match(pattern, fact: String)
do
  let b := match_atom(atom_of(pattern), atom_of(fact), create Bindings.empty)
  if b = nil then
    console.print_line(pattern + " against " + fact + ": no match")
  else
    console.print_line(pattern + " against " + fact + ": " + b.text)
  end
end

try_match("follows(X, Y)", "follows(ann, bo)")
try_match("follows(X, bo)", "follows(ann, bo)")
try_match("follows(X, bo)", "follows(ann, cy)")
try_match("follows(X, X)", "follows(ann, bo)")
try_match("follows(X, X)", "follows(bo, bo)")
try_match("follows(X, Y)", "sees(ann, bo)")
console.print_line(holds("<", create Int.make(480), create Int.make(1000)).to_string)
console.print_line(holds("!=", create Symbol.make("bob"), create Symbol.make("bob")).to_string)
\end{tryit}
\tryoutput{try_matching}

The try program reads its patterns with the parser, by asking it to parse a
question, which is the simplest way to get an atom from text.
